\chapter{Phát triển Hệ E-Commerce Microservices}

\section{Mục tiêu chương}

Chương này trình bày quá trình phân tích và thiết kế một hệ thống thương mại điện tử theo kiến trúc Microservices. Khác với việc chỉ mô tả lý thuyết, nội dung chương được xây dựng theo hướng có thể sử dụng làm tài liệu thiết kế ban đầu để triển khai hệ thống thực tế. Vì vậy, ngoài việc xác định các service chính, chương còn mô tả rõ trách nhiệm của từng service, dữ liệu quản lý, API contract, luồng nghiệp vụ, trạng thái xử lý, nguyên tắc giao tiếp giữa các service và thiết kế cơ sở dữ liệu.

Hệ thống được thiết kế theo hướng một nền tảng ecommerce cơ bản, cho phép người dùng xem sản phẩm, thêm sản phẩm vào giỏ hàng, đặt hàng, thanh toán và theo dõi giao hàng. Bên cạnh đó, hệ thống cũng hỗ trợ các chức năng quản trị như quản lý sản phẩm, quản lý người dùng, quản lý đơn hàng và cập nhật trạng thái vận hành.

Mục tiêu chính của chương gồm:

\begin{itemize}
  \item Xác định yêu cầu chức năng và phi chức năng của hệ thống ecommerce.
  \item Phân rã hệ thống thành các microservice dựa trên Domain Driven Design.
  \item Thiết kế trách nhiệm, dữ liệu và API cho từng service.
  \item Mô tả luồng checkout có thể triển khai thực tế.
  \item Đảm bảo các vấn đề quan trọng như giá tiền, tồn kho, thanh toán, giao hàng và trạng thái đơn hàng được xử lý hợp lý.
  \item Đề xuất database schema và cấu hình triển khai ban đầu.
\end{itemize}

\section{Xác định yêu cầu hệ thống}

\subsection{Functional Requirements}

Hệ thống ecommerce cần đáp ứng các yêu cầu chức năng sau:

\begin{itemize}
  \item \textbf{Quản lý người dùng}: Người dùng có thể đăng ký, đăng nhập, cập nhật thông tin cá nhân và sử dụng chức năng phù hợp với vai trò của mình. Các vai trò chính gồm Admin, Staff và Customer.

  \item \textbf{Xác thực và phân quyền}: Hệ thống sử dụng JWT để xác thực người dùng. Các API quan trọng cần kiểm tra token và quyền truy cập trước khi xử lý request.

  \item \textbf{Quản lý sản phẩm}: Admin hoặc Staff có thể thêm, sửa, xóa và cập nhật thông tin sản phẩm. Customer có thể xem danh sách sản phẩm, xem chi tiết sản phẩm và tìm kiếm sản phẩm.

  \item \textbf{Quản lý biến thể sản phẩm}: Một sản phẩm có thể có nhiều biến thể khác nhau, ví dụ cùng một mẫu áo có nhiều size, màu sắc; cùng một điện thoại có nhiều màu và dung lượng bộ nhớ. Vì vậy hệ thống cần hỗ trợ Product Variant hoặc SKU.

  \item \textbf{Quản lý tồn kho}: Hệ thống cần kiểm tra số lượng tồn kho trước khi tạo đơn hàng. Trong phạm vi thiết kế hiện tại, tồn kho được quản lý trong Product Service. Khi hệ thống mở rộng, phần tồn kho có thể được tách thành Inventory Service riêng.

  \item \textbf{Quản lý giỏ hàng}: Customer có thể thêm sản phẩm vào giỏ hàng, cập nhật số lượng, xóa sản phẩm khỏi giỏ và xem tổng giá trị tạm tính.

  \item \textbf{Checkout}: Khi Customer xác nhận mua hàng, hệ thống kiểm tra giỏ hàng, xác thực lại giá, kiểm tra tồn kho, giữ hàng tạm thời, tạo đơn hàng và tạo yêu cầu thanh toán.

  \item \textbf{Quản lý đơn hàng}: Hệ thống lưu đơn hàng, chi tiết đơn hàng, trạng thái đơn hàng và lịch sử xử lý đơn hàng. Customer có thể xem đơn hàng của mình, Staff hoặc Admin có thể xử lý đơn hàng.

  \item \textbf{Xử lý thanh toán}: Hệ thống tạo giao dịch thanh toán cho đơn hàng. Payment Service có thể mô phỏng thanh toán hoặc tích hợp với cổng thanh toán bên ngoài như VNPay, Momo, ZaloPay hoặc PayPal trong tương lai.

  \item \textbf{Quản lý giao hàng}: Sau khi thanh toán thành công, hệ thống tạo thông tin giao hàng và cập nhật trạng thái vận chuyển.

  \item \textbf{Thông báo}: Hệ thống có thể gửi email hoặc thông báo cho người dùng khi đơn hàng được tạo, thanh toán thành công hoặc giao hàng hoàn tất.
\end{itemize}

\subsection{Non-functional Requirements}

Các yêu cầu phi chức năng đóng vai trò quan trọng vì ecommerce là hệ thống có nhiều tương tác, nhiều dữ liệu và cần độ ổn định cao.

\begin{itemize}
  \item \textbf{Scalability}: Hệ thống cần có khả năng mở rộng từng service độc lập. Ví dụ, khi lượng người dùng xem sản phẩm tăng cao, chỉ cần mở rộng Product Service thay vì toàn bộ hệ thống.

  \item \textbf{Availability}: Các chức năng chính như xem sản phẩm, đặt hàng và thanh toán cần có tính sẵn sàng cao. Nếu một service phụ như Notification Service gặp lỗi, hệ thống vẫn cần cho phép người dùng tiếp tục đặt hàng.

  \item \textbf{Security}: Hệ thống cần xác thực bằng JWT, mã hóa mật khẩu, phân quyền theo vai trò và giới hạn quyền truy cập API.

  \item \textbf{Maintainability}: Mỗi service cần có trách nhiệm rõ ràng, codebase riêng và database riêng để dễ bảo trì và nâng cấp.

  \item \textbf{Data Consistency}: Vì dữ liệu được chia giữa nhiều service, hệ thống cần chấp nhận mô hình nhất quán sau cùng trong một số luồng nghiệp vụ. Các nghiệp vụ quan trọng như Order, Payment và Stock cần được thiết kế cẩn thận để tránh sai lệch dữ liệu.

  \item \textbf{Observability}: Hệ thống cần logging, monitoring và tracing để theo dõi request đi qua nhiều service.

  \item \textbf{Performance}: Các API đọc nhiều như danh sách sản phẩm, chi tiết sản phẩm và giỏ hàng cần phản hồi nhanh. Có thể sử dụng cache cho dữ liệu ít thay đổi.
\end{itemize}

\section{Phân rã hệ thống theo Domain Driven Design}

\subsection{Bounded Context}

Theo Domain Driven Design, hệ thống nên được chia theo ranh giới nghiệp vụ thay vì chia theo tầng kỹ thuật. Với hệ thống ecommerce, các Bounded Context chính gồm:

\begin{itemize}
  \item \textbf{User Context}: Quản lý tài khoản, đăng nhập, phân quyền và hồ sơ người dùng.
  \item \textbf{Product Context}: Quản lý sản phẩm, danh mục, biến thể sản phẩm và tồn kho cơ bản.
  \item \textbf{Cart Context}: Quản lý giỏ hàng tạm thời của khách hàng.
  \item \textbf{Order Context}: Quản lý đơn hàng, chi tiết đơn hàng và vòng đời đơn hàng.
  \item \textbf{Payment Context}: Quản lý thanh toán và trạng thái giao dịch.
  \item \textbf{Shipping Context}: Quản lý vận chuyển và trạng thái giao hàng.
  \item \textbf{Notification Context}: Quản lý gửi email, SMS hoặc thông báo hệ thống.
\end{itemize}

Mỗi Bounded Context được ánh xạ thành một microservice độc lập như sau:

\begin{longtable}{|p{4cm}|p{4cm}|p{6cm}|}
\hline
\textbf{Bounded Context} & \textbf{Service} & \textbf{Trách nhiệm chính} \\
\hline
User Context & \texttt{user-service} & Quản lý người dùng, xác thực, phân quyền \\
\hline
Product Context & \texttt{product-service} & Quản lý sản phẩm, danh mục, SKU, tồn kho cơ bản \\
\hline
Cart Context & \texttt{cart-service} & Quản lý giỏ hàng tạm thời của customer \\
\hline
Order Context & \texttt{order-service} & Điều phối checkout, tạo đơn hàng, quản lý trạng thái đơn hàng \\
\hline
Payment Context & \texttt{payment-service} & Tạo giao dịch, xử lý callback/webhook thanh toán \\
\hline
Shipping Context & \texttt{shipping-service} & Tạo vận đơn, cập nhật trạng thái vận chuyển \\
\hline
Notification Context & \texttt{notification-service} & Gửi email, SMS hoặc thông báo cho người dùng \\
\hline
\end{longtable}

\subsection{Nguyên tắc thiết kế}

Khi thiết kế hệ thống theo Microservices, cần tuân thủ các nguyên tắc sau:

\begin{itemize}
  \item \textbf{Database per Service}: Mỗi service có database riêng. Service khác không được truy cập trực tiếp vào database nội bộ của service đó.

  \item \textbf{API Contract rõ ràng}: Các service giao tiếp thông qua API hoặc message. Dữ liệu request và response cần được định nghĩa rõ để tránh phụ thuộc ngầm.

  \item \textbf{Không dùng Foreign Key giữa các database}: Các ID như \texttt{user\_id}, \texttt{product\_variant\_id}, \texttt{order\_id} chỉ là tham chiếu logic giữa các service.

  \item \textbf{Lưu dữ liệu snapshot}: Order Service cần lưu lại tên sản phẩm, giá tại thời điểm mua và số lượng. Điều này đảm bảo đơn hàng cũ không bị thay đổi khi Product Service cập nhật thông tin sản phẩm.

  \item \textbf{Thiết kế cho lỗi}: Hệ thống phân tán có thể gặp timeout, mất kết nối hoặc service tạm thời không phản hồi. Vì vậy cần có retry, timeout, fallback và logging.

  \item \textbf{Không chia service quá nhỏ ở giai đoạn đầu}: Trong phạm vi hệ thống hiện tại, Product Service tạm thời quản lý tồn kho. Khi hệ thống lớn hơn, có thể tách tồn kho thành Inventory Service riêng.
\end{itemize}

\section{Kiến trúc tổng thể hệ thống}

\subsection{Mô hình kiến trúc}

Kiến trúc tổng thể của hệ thống được mô tả như sau:

\begin{verbatim}
Client / Frontend
        |
        v
   API Gateway
        |
        |----> user-service
        |----> product-service
        |----> cart-service
        |----> order-service
        |           |
        |           |----> payment-service
        |           |----> shipping-service
        |           |----> notification-service
        |
        |----> notification-service
\end{verbatim}

API Gateway là điểm vào duy nhất của hệ thống. Client không gọi trực tiếp các service nội bộ mà gửi request đến Gateway. Gateway chịu trách nhiệm định tuyến request, kiểm tra token cơ bản, giới hạn tần suất truy cập và chuyển request đến service phù hợp.

\subsection{Danh sách service và port đề xuất}

\begin{longtable}{|p{4cm}|p{4cm}|p{5cm}|}
\hline
\textbf{Service} & \textbf{Port đề xuất} & \textbf{Database} \\
\hline
API Gateway & 8080 & Không bắt buộc có database \\
\hline
user-service & 8001 & user\_db \\
\hline
product-service & 8002 & product\_db \\
\hline
cart-service & 8003 & cart\_db hoặc Redis \\
\hline
order-service & 8004 & order\_db \\
\hline
payment-service & 8005 & payment\_db \\
\hline
shipping-service & 8006 & shipping\_db \\
\hline
notification-service & 8007 & notification\_db hoặc message queue \\
\hline
\end{longtable}

\section{Thiết kế User Service}

\subsection{Trách nhiệm}

\texttt{user-service} chịu trách nhiệm quản lý người dùng, xác thực và phân quyền. Service này không xử lý đơn hàng, giỏ hàng hay thanh toán. Các service khác chỉ sử dụng \texttt{user\_id} như một tham chiếu logic.

\subsection{Vai trò người dùng}

\begin{longtable}{|p{3cm}|p{10cm}|}
\hline
\textbf{Vai trò} & \textbf{Quyền hạn} \\
\hline
Admin & Quản lý toàn bộ hệ thống, bao gồm người dùng, sản phẩm, đơn hàng và cấu hình. \\
\hline
Staff & Hỗ trợ vận hành đơn hàng, cập nhật vận chuyển, xử lý yêu cầu khách hàng. \\
\hline
Customer & Xem sản phẩm, thêm vào giỏ hàng, đặt hàng, thanh toán và theo dõi đơn hàng. \\
\hline
\end{longtable}

\subsection{Django Model}

\begin{lstlisting}[style=code,language=Python]
from django.db import models
from django.contrib.auth.models import AbstractUser


class User(AbstractUser):
    ROLE_CHOICES = (
        ("admin", "Admin"),
        ("staff", "Staff"),
        ("customer", "Customer"),
    )

    role = models.CharField(
        max_length=20,
        choices=ROLE_CHOICES,
        default="customer"
    )
    phone = models.CharField(max_length=20, null=True, blank=True)
    address = models.TextField(null=True, blank=True)
    is_active = models.BooleanField(default=True)
    created_at = models.DateTimeField(auto_now_add=True)
    updated_at = models.DateTimeField(auto_now=True)
\end{lstlisting}

\subsection{API Contract}

\begin{lstlisting}[style=code]
POST /api/auth/register/
Request:
{
  "username": "customer01",
  "email": "customer01@example.com",
  "password": "123456"
}

Response:
{
  "id": 1,
  "username": "customer01",
  "email": "customer01@example.com",
  "role": "customer"
}
\end{lstlisting}

\begin{lstlisting}[style=code]
POST /api/auth/login/
Request:
{
  "username": "customer01",
  "password": "123456"
}

Response:
{
  "access_token": "jwt_access_token",
  "refresh_token": "jwt_refresh_token",
  "user": {
    "id": 1,
    "username": "customer01",
    "role": "customer"
  }
}
\end{lstlisting}

Các API chính của User Service:

\begin{lstlisting}[style=code]
POST /api/auth/register/
POST /api/auth/login/
POST /api/auth/logout/
GET  /api/users/me/
PUT  /api/users/me/
GET  /api/users/
GET  /api/users/{id}/
PUT  /api/users/{id}/role/
\end{lstlisting}

\section{Thiết kế Product Service}

\subsection{Trách nhiệm}

\texttt{product-service} quản lý thông tin sản phẩm, danh mục, biến thể sản phẩm và tồn kho cơ bản. Trong thiết kế này, Product Service tạm thời chịu trách nhiệm quản lý tồn kho để giảm số lượng service ban đầu. Khi hệ thống mở rộng, phần tồn kho có thể được tách thành \texttt{inventory-service} riêng.

\subsection{Product, Variant và SKU}

Trong ecommerce thực tế, một sản phẩm có thể có nhiều biến thể. Vì vậy không nên chỉ lưu mỗi bảng \texttt{Product}. Hệ thống cần thêm \texttt{ProductVariant} để biểu diễn từng phiên bản có thể bán được.

Ví dụ:

\begin{verbatim}
Product: iPhone 15
  Variant 1: iPhone 15 - Black - 128GB
  Variant 2: iPhone 15 - Blue - 256GB

Product: Áo thun basic
  Variant 1: Áo thun basic - Black - Size M
  Variant 2: Áo thun basic - White - Size L
\end{verbatim}

\subsection{Django Model}

\begin{lstlisting}[style=code,language=Python]
from django.db import models


class Category(models.Model):
    name = models.CharField(max_length=100)
    description = models.TextField(null=True, blank=True)
    is_active = models.BooleanField(default=True)


class Product(models.Model):
    PRODUCT_TYPES = (
        ("book", "Book"),
        ("electronics", "Electronics"),
        ("fashion", "Fashion"),
    )

    category = models.ForeignKey(
        Category,
        on_delete=models.CASCADE,
        related_name="products"
    )
    name = models.CharField(max_length=255)
    description = models.TextField(null=True, blank=True)
    product_type = models.CharField(max_length=50, choices=PRODUCT_TYPES)
    brand = models.CharField(max_length=100, null=True, blank=True)
    is_active = models.BooleanField(default=True)
    created_at = models.DateTimeField(auto_now_add=True)
    updated_at = models.DateTimeField(auto_now=True)


class ProductVariant(models.Model):
    product = models.ForeignKey(
        Product,
        on_delete=models.CASCADE,
        related_name="variants"
    )
    sku = models.CharField(max_length=100, unique=True)
    variant_name = models.CharField(max_length=255)
    price = models.DecimalField(max_digits=12, decimal_places=2)
    stock_quantity = models.IntegerField(default=0)
    reserved_quantity = models.IntegerField(default=0)
    attributes = models.JSONField(null=True, blank=True)
    is_active = models.BooleanField(default=True)

    def available_stock(self):
        return self.stock_quantity - self.reserved_quantity
\end{lstlisting}

Trường \texttt{attributes} có thể lưu các thuộc tính linh hoạt theo từng loại sản phẩm. Ví dụ:

\begin{lstlisting}[style=code]
{
  "color": "black",
  "size": "M"
}
\end{lstlisting}

hoặc:

\begin{lstlisting}[style=code]
{
  "color": "blue",
  "storage": "256GB",
  "warranty_months": 12
}
\end{lstlisting}

\subsection{Logic tồn kho}

Để tránh tình trạng bán quá số lượng tồn kho, hệ thống cần có cơ chế giữ hàng tạm thời khi người dùng checkout.

Luồng xử lý tồn kho đề xuất:

\begin{enumerate}
  \item Khi Customer checkout, Order Service yêu cầu Product Service kiểm tra tồn kho.
  \item Nếu đủ hàng, Product Service tăng \texttt{reserved\_quantity} để giữ hàng tạm thời.
  \item Nếu thanh toán thành công, Product Service trừ \texttt{stock\_quantity} và giảm \texttt{reserved\_quantity}.
  \item Nếu thanh toán thất bại hoặc hết hạn thanh toán, Product Service giảm \texttt{reserved\_quantity} để trả hàng về kho.
\end{enumerate}

Cách xử lý này giúp hạn chế lỗi overselling khi nhiều người dùng cùng mua một sản phẩm có số lượng tồn kho thấp.

\subsection{API Contract}

\begin{lstlisting}[style=code]
GET /api/products/
Response:
[
  {
    "id": 1,
    "name": "iPhone 15",
    "product_type": "electronics",
    "brand": "Apple",
    "variants": [
      {
        "id": 10,
        "sku": "IPHONE15-BLACK-128",
        "variant_name": "Black 128GB",
        "price": "19990000.00",
        "available_stock": 20
      }
    ]
  }
]
\end{lstlisting}

\begin{lstlisting}[style=code]
POST /api/products/variants/{id}/reserve/
Request:
{
  "quantity": 2,
  "order_id": 1001
}

Response:
{
  "variant_id": 10,
  "reserved_quantity": 2,
  "available_stock": 18
}
\end{lstlisting}

Các API chính của Product Service:

\begin{lstlisting}[style=code]
GET    /api/products/
GET    /api/products/{id}/
POST   /api/products/
PUT    /api/products/{id}/
DELETE /api/products/{id}/

GET    /api/products/variants/{id}/
POST   /api/products/variants/
PUT    /api/products/variants/{id}/

POST   /api/products/variants/{id}/reserve/
POST   /api/products/variants/{id}/confirm-stock/
POST   /api/products/variants/{id}/release-stock/
\end{lstlisting}

\section{Thiết kế Cart Service}

\subsection{Trách nhiệm}

\texttt{cart-service} quản lý giỏ hàng tạm thời của người dùng. Cart Service không phải nơi quyết định giá cuối cùng của đơn hàng. Giá trong giỏ hàng chỉ dùng để hiển thị tạm thời. Khi checkout, Order Service cần kiểm tra lại giá và tồn kho từ Product Service.

\subsection{Django Model}

\begin{lstlisting}[style=code,language=Python]
from django.db import models


class Cart(models.Model):
    user_id = models.IntegerField()
    created_at = models.DateTimeField(auto_now_add=True)
    updated_at = models.DateTimeField(auto_now=True)


class CartItem(models.Model):
    cart = models.ForeignKey(
        Cart,
        on_delete=models.CASCADE,
        related_name="items"
    )
    product_id = models.IntegerField()
    product_variant_id = models.IntegerField()
    product_name = models.CharField(max_length=255)
    variant_name = models.CharField(max_length=255)
    unit_price = models.DecimalField(max_digits=12, decimal_places=2)
    quantity = models.IntegerField(default=1)

    def subtotal(self):
        return self.unit_price * self.quantity
\end{lstlisting}

\subsection{API Contract}

\begin{lstlisting}[style=code]
POST /api/cart/items/
Request:
{
  "product_variant_id": 10,
  "quantity": 2
}

Response:
{
  "cart_id": 1,
  "items": [
    {
      "product_variant_id": 10,
      "product_name": "iPhone 15",
      "variant_name": "Black 128GB",
      "unit_price": "19990000.00",
      "quantity": 2,
      "subtotal": "39980000.00"
    }
  ]
}
\end{lstlisting}

Các API chính của Cart Service:

\begin{lstlisting}[style=code]
GET    /api/cart/
POST   /api/cart/items/
PUT    /api/cart/items/{item_id}/
DELETE /api/cart/items/{item_id}/
DELETE /api/cart/clear/
\end{lstlisting}

\section{Thiết kế Order Service}

\subsection{Trách nhiệm}

\texttt{order-service} là service trung tâm trong luồng checkout. Service này không chỉ lưu đơn hàng mà còn điều phối quá trình đặt hàng, kiểm tra giỏ hàng, xác thực lại giá sản phẩm, yêu cầu giữ hàng, tạo thanh toán và tạo giao hàng sau khi thanh toán thành công.

\subsection{Order State Machine}

Một đơn hàng cần có vòng đời rõ ràng. Các trạng thái đề xuất:

\begin{longtable}{|p{4cm}|p{9cm}|}
\hline
\textbf{Trạng thái} & \textbf{Ý nghĩa} \\
\hline
created & Đơn hàng vừa được tạo trong hệ thống. \\
\hline
pending\_payment & Đơn hàng đang chờ thanh toán. \\
\hline
paid & Đơn hàng đã thanh toán thành công. \\
\hline
preparing & Đơn hàng đang được chuẩn bị. \\
\hline
shipping & Đơn hàng đang được giao. \\
\hline
completed & Đơn hàng đã hoàn tất. \\
\hline
cancelled & Đơn hàng đã bị hủy. \\
\hline
failed & Đơn hàng thất bại do thanh toán lỗi hoặc lỗi xử lý. \\
\hline
refunded & Đơn hàng đã được hoàn tiền. \\
\hline
\end{longtable}

Các chuyển trạng thái hợp lệ:

\begin{verbatim}
created -> pending_payment
pending_payment -> paid
pending_payment -> failed
pending_payment -> cancelled
paid -> preparing
preparing -> shipping
shipping -> completed
paid -> refunded
\end{verbatim}

\subsection{Django Model}

\begin{lstlisting}[style=code,language=Python]
from django.db import models


class Order(models.Model):
    STATUS_CHOICES = (
        ("created", "Created"),
        ("pending_payment", "Pending Payment"),
        ("paid", "Paid"),
        ("preparing", "Preparing"),
        ("shipping", "Shipping"),
        ("completed", "Completed"),
        ("cancelled", "Cancelled"),
        ("failed", "Failed"),
        ("refunded", "Refunded"),
    )

    user_id = models.IntegerField()
    total_amount = models.DecimalField(max_digits=12, decimal_places=2)
    status = models.CharField(
        max_length=30,
        choices=STATUS_CHOICES,
        default="created"
    )
    shipping_address = models.TextField()
    payment_method = models.CharField(max_length=50)
    created_at = models.DateTimeField(auto_now_add=True)
    updated_at = models.DateTimeField(auto_now=True)


class OrderItem(models.Model):
    order = models.ForeignKey(
        Order,
        on_delete=models.CASCADE,
        related_name="items"
    )
    product_id = models.IntegerField()
    product_variant_id = models.IntegerField()
    product_name = models.CharField(max_length=255)
    variant_name = models.CharField(max_length=255)
    unit_price = models.DecimalField(max_digits=12, decimal_places=2)
    quantity = models.IntegerField()
    subtotal = models.DecimalField(max_digits=12, decimal_places=2)
\end{lstlisting}

\subsection{Checkout Flow}

Luồng checkout được thiết kế như sau:

\begin{enumerate}
  \item Customer gửi yêu cầu checkout đến Order Service.
  \item Order Service gọi Cart Service để lấy danh sách sản phẩm trong giỏ hàng.
  \item Order Service gọi Product Service để xác thực lại sản phẩm, giá hiện tại và tồn kho.
  \item Nếu sản phẩm không tồn tại hoặc không đủ hàng, hệ thống trả lỗi.
  \item Nếu hợp lệ, Order Service yêu cầu Product Service giữ hàng tạm thời.
  \item Order Service tạo Order với trạng thái \texttt{pending\_payment}.
  \item Order Service tạo OrderItem bằng dữ liệu snapshot tại thời điểm mua.
  \item Order Service gọi Payment Service để tạo giao dịch thanh toán.
  \item Payment Service trả về \texttt{payment\_url} hoặc \texttt{transaction\_code}.
  \item Customer thực hiện thanh toán.
  \item Payment Service nhận callback/webhook và cập nhật trạng thái thanh toán.
  \item Nếu thanh toán thành công, Order Service cập nhật đơn hàng sang \texttt{paid}, xác nhận trừ kho và yêu cầu Shipping Service tạo vận đơn.
  \item Nếu thanh toán thất bại, Order Service cập nhật đơn hàng sang \texttt{failed} và yêu cầu Product Service release stock.
\end{enumerate}

\subsection{API Contract}

\begin{lstlisting}[style=code]
POST /api/orders/checkout/
Request:
{
  "shipping_address": "Ha Noi, Viet Nam",
  "payment_method": "momo"
}

Response:
{
  "order_id": 1001,
  "status": "pending_payment",
  "total_amount": "39980000.00",
  "payment": {
    "payment_id": 5001,
    "payment_url": "https://payment-gateway.example/pay/5001"
  }
}
\end{lstlisting}

Các API chính của Order Service:

\begin{lstlisting}[style=code]
POST /api/orders/checkout/
GET  /api/orders/
GET  /api/orders/{id}/
PUT  /api/orders/{id}/status/
POST /api/orders/{id}/cancel/
POST /api/orders/{id}/payment-success/
POST /api/orders/{id}/payment-failed/
\end{lstlisting}

\section{Thiết kế Payment Service}

\subsection{Trách nhiệm}

\texttt{payment-service} chịu trách nhiệm tạo giao dịch thanh toán, lưu trạng thái thanh toán và xử lý callback hoặc webhook từ cổng thanh toán. Trong hệ thống thực tế, thanh toán thường không hoàn tất ngay tại thời điểm tạo request. Vì vậy, Payment Service cần hỗ trợ cơ chế bất đồng bộ.

\subsection{Payment State Machine}

\begin{longtable}{|p{4cm}|p{9cm}|}
\hline
\textbf{Trạng thái} & \textbf{Ý nghĩa} \\
\hline
pending & Giao dịch đã được tạo nhưng chưa xử lý. \\
\hline
processing & Giao dịch đang được xử lý bởi cổng thanh toán. \\
\hline
success & Thanh toán thành công. \\
\hline
failed & Thanh toán thất bại. \\
\hline
cancelled & Người dùng hủy giao dịch. \\
\hline
refunded & Giao dịch đã được hoàn tiền. \\
\hline
\end{longtable}

\subsection{Django Model}

\begin{lstlisting}[style=code,language=Python]
from django.db import models


class Payment(models.Model):
    STATUS_CHOICES = (
        ("pending", "Pending"),
        ("processing", "Processing"),
        ("success", "Success"),
        ("failed", "Failed"),
        ("cancelled", "Cancelled"),
        ("refunded", "Refunded"),
    )

    order_id = models.IntegerField()
    user_id = models.IntegerField()
    amount = models.DecimalField(max_digits=12, decimal_places=2)
    payment_method = models.CharField(max_length=50)
    status = models.CharField(
        max_length=20,
        choices=STATUS_CHOICES,
        default="pending"
    )
    transaction_code = models.CharField(max_length=100, null=True, blank=True)
    payment_url = models.TextField(null=True, blank=True)
    gateway_response = models.JSONField(null=True, blank=True)
    created_at = models.DateTimeField(auto_now_add=True)
    updated_at = models.DateTimeField(auto_now=True)
\end{lstlisting}

\subsection{Payment Flow}

\begin{enumerate}
  \item Order Service gửi yêu cầu tạo thanh toán.
  \item Payment Service tạo bản ghi Payment với trạng thái \texttt{pending}.
  \item Payment Service tạo \texttt{payment\_url} hoặc \texttt{transaction\_code}.
  \item Customer được chuyển đến trang thanh toán.
  \item Sau khi thanh toán, cổng thanh toán gửi webhook về Payment Service.
  \item Payment Service xác thực webhook và cập nhật trạng thái.
  \item Payment Service thông báo kết quả cho Order Service.
\end{enumerate}

\subsection{API Contract}

\begin{lstlisting}[style=code]
POST /api/payments/
Request:
{
  "order_id": 1001,
  "user_id": 1,
  "amount": "39980000.00",
  "payment_method": "momo"
}

Response:
{
  "payment_id": 5001,
  "order_id": 1001,
  "status": "pending",
  "payment_url": "https://payment-gateway.example/pay/5001"
}
\end{lstlisting}

\begin{lstlisting}[style=code]
POST /api/payments/webhook/
Request:
{
  "transaction_code": "TXN123456",
  "order_id": 1001,
  "status": "success",
  "amount": "39980000.00"
}

Response:
{
  "message": "Payment status updated"
}
\end{lstlisting}

Các API chính của Payment Service:

\begin{lstlisting}[style=code]
POST /api/payments/
GET  /api/payments/{id}/
GET  /api/payments/order/{order_id}/
POST /api/payments/webhook/
POST /api/payments/{id}/refund/
\end{lstlisting}

\section{Thiết kế Shipping Service}

\subsection{Trách nhiệm}

\texttt{shipping-service} chịu trách nhiệm tạo và quản lý thông tin vận chuyển sau khi đơn hàng được thanh toán thành công. Service này không xử lý thanh toán và không tính lại giá đơn hàng.

\subsection{Shipping State Machine}

\begin{longtable}{|p{4cm}|p{9cm}|}
\hline
\textbf{Trạng thái} & \textbf{Ý nghĩa} \\
\hline
processing & Đơn hàng đang được chuẩn bị để giao. \\
\hline
shipping & Đơn hàng đang được vận chuyển. \\
\hline
delivered & Đơn hàng đã giao thành công. \\
\hline
failed & Giao hàng thất bại. \\
\hline
returned & Đơn hàng bị hoàn trả. \\
\hline
\end{longtable}

\subsection{Django Model}

\begin{lstlisting}[style=code,language=Python]
from django.db import models


class Shipment(models.Model):
    STATUS_CHOICES = (
        ("processing", "Processing"),
        ("shipping", "Shipping"),
        ("delivered", "Delivered"),
        ("failed", "Failed"),
        ("returned", "Returned"),
    )

    order_id = models.IntegerField()
    user_id = models.IntegerField()
    receiver_name = models.CharField(max_length=255)
    receiver_phone = models.CharField(max_length=20)
    address = models.TextField()
    carrier = models.CharField(max_length=100, null=True, blank=True)
    tracking_code = models.CharField(max_length=100, null=True, blank=True)
    status = models.CharField(
        max_length=30,
        choices=STATUS_CHOICES,
        default="processing"
    )
    created_at = models.DateTimeField(auto_now_add=True)
    updated_at = models.DateTimeField(auto_now=True)
\end{lstlisting}

\subsection{API Contract}

\begin{lstlisting}[style=code]
POST /api/shipments/
Request:
{
  "order_id": 1001,
  "user_id": 1,
  "receiver_name": "Tran Xuan Thanh",
  "receiver_phone": "0912345678",
  "address": "Ha Noi, Viet Nam"
}

Response:
{
  "shipment_id": 7001,
  "order_id": 1001,
  "status": "processing",
  "tracking_code": "SHIP1001"
}
\end{lstlisting}

Các API chính của Shipping Service:

\begin{lstlisting}[style=code]
POST /api/shipments/
GET  /api/shipments/{id}/
GET  /api/shipments/order/{order_id}/
PUT  /api/shipments/{id}/status/
\end{lstlisting}

\section{Thiết kế Notification Service}

\subsection{Trách nhiệm}

\texttt{notification-service} chịu trách nhiệm gửi thông báo cho người dùng. Service này nên được thiết kế bất đồng bộ để không làm chậm luồng nghiệp vụ chính. Ví dụ, nếu gửi email thất bại, đơn hàng vẫn không nên bị hủy.

Các loại thông báo gồm:

\begin{itemize}
  \item Thông báo tạo đơn hàng thành công.
  \item Thông báo thanh toán thành công hoặc thất bại.
  \item Thông báo đơn hàng đang được giao.
  \item Thông báo đơn hàng hoàn tất.
\end{itemize}

\subsection{API chính}

\begin{lstlisting}[style=code]
POST /api/notifications/email/
POST /api/notifications/sms/
GET  /api/notifications/user/{user_id}/
\end{lstlisting}

Trong phiên bản nâng cao, Notification Service nên nhận event từ message broker thay vì được gọi trực tiếp bằng REST API.

\section{Giao tiếp giữa các service}

\subsection{Giao tiếp đồng bộ}

REST API được sử dụng cho các thao tác cần phản hồi ngay, ví dụ:

\begin{longtable}{|p{4cm}|p{4cm}|p{6cm}|}
\hline
\textbf{Service gọi} & \textbf{Service được gọi} & \textbf{Mục đích} \\
\hline
Cart Service & Product Service & Lấy thông tin sản phẩm khi thêm vào giỏ \\
\hline
Order Service & Cart Service & Lấy dữ liệu giỏ hàng khi checkout \\
\hline
Order Service & Product Service & Kiểm tra giá, tồn kho và giữ hàng \\
\hline
Order Service & Payment Service & Tạo giao dịch thanh toán \\
\hline
Order Service & Shipping Service & Tạo vận đơn sau khi thanh toán thành công \\
\hline
\end{longtable}

\subsection{Giao tiếp bất đồng bộ}

Một số thao tác không cần xử lý ngay lập tức có thể dùng message broker như RabbitMQ hoặc Kafka.

Ví dụ event:

\begin{lstlisting}[style=code]
OrderCreated
PaymentSucceeded
PaymentFailed
ShipmentCreated
OrderCompleted
\end{lstlisting}

Giao tiếp bất đồng bộ giúp giảm coupling giữa các service và tăng khả năng chịu lỗi. Ví dụ, khi thanh toán thành công, Payment Service có thể phát event \texttt{PaymentSucceeded}. Order Service nhận event này để cập nhật trạng thái đơn hàng, Shipping Service tạo vận đơn và Notification Service gửi email xác nhận.

\section{Thiết kế Database cho từng Service}

\subsection{Nguyên tắc chung}

Các database được thiết kế theo nguyên tắc database-per-service. Những ID liên service chỉ là tham chiếu logic, không phải foreign key vật lý giữa các database.

Ngoài ra, các trường tiền tệ như \texttt{price}, \texttt{amount}, \texttt{total\_amount} không sử dụng kiểu \texttt{FLOAT}. Thay vào đó, hệ thống sử dụng \texttt{DECIMAL(12,2)} để tránh sai số khi tính toán tiền.

\subsection{User Service Database}

\begin{lstlisting}[style=code,language=SQL]
CREATE TABLE users (
    id SERIAL PRIMARY KEY,
    username VARCHAR(100) UNIQUE NOT NULL,
    email VARCHAR(255) UNIQUE NOT NULL,
    password VARCHAR(255) NOT NULL,
    role VARCHAR(20) DEFAULT 'customer',
    phone VARCHAR(20),
    address TEXT,
    is_active BOOLEAN DEFAULT TRUE,
    created_at TIMESTAMP DEFAULT CURRENT_TIMESTAMP,
    updated_at TIMESTAMP DEFAULT CURRENT_TIMESTAMP
);
\end{lstlisting}

\subsection{Product Service Database}

\begin{lstlisting}[style=code,language=SQL]
CREATE TABLE categories (
    id SERIAL PRIMARY KEY,
    name VARCHAR(100) NOT NULL,
    description TEXT,
    is_active BOOLEAN DEFAULT TRUE
);

CREATE TABLE products (
    id SERIAL PRIMARY KEY,
    category_id INT REFERENCES categories(id),
    name VARCHAR(255) NOT NULL,
    description TEXT,
    product_type VARCHAR(50),
    brand VARCHAR(100),
    is_active BOOLEAN DEFAULT TRUE,
    created_at TIMESTAMP DEFAULT CURRENT_TIMESTAMP,
    updated_at TIMESTAMP DEFAULT CURRENT_TIMESTAMP
);

CREATE TABLE product_variants (
    id SERIAL PRIMARY KEY,
    product_id INT REFERENCES products(id),
    sku VARCHAR(100) UNIQUE NOT NULL,
    variant_name VARCHAR(255),
    price DECIMAL(12, 2) NOT NULL,
    stock_quantity INT DEFAULT 0,
    reserved_quantity INT DEFAULT 0,
    attributes JSONB,
    is_active BOOLEAN DEFAULT TRUE
);

CREATE INDEX idx_product_variants_sku
ON product_variants(sku);
\end{lstlisting}

\subsection{Cart Service Database}

\begin{lstlisting}[style=code,language=SQL]
CREATE TABLE carts (
    id SERIAL PRIMARY KEY,
    user_id INT NOT NULL,
    created_at TIMESTAMP DEFAULT CURRENT_TIMESTAMP,
    updated_at TIMESTAMP DEFAULT CURRENT_TIMESTAMP
);

CREATE TABLE cart_items (
    id SERIAL PRIMARY KEY,
    cart_id INT REFERENCES carts(id),
    product_id INT NOT NULL,
    product_variant_id INT NOT NULL,
    product_name VARCHAR(255),
    variant_name VARCHAR(255),
    unit_price DECIMAL(12, 2) NOT NULL,
    quantity INT NOT NULL
);

CREATE INDEX idx_carts_user_id
ON carts(user_id);
\end{lstlisting}

\subsection{Order Service Database}

\begin{lstlisting}[style=code,language=SQL]
CREATE TABLE orders (
    id SERIAL PRIMARY KEY,
    user_id INT NOT NULL,
    total_amount DECIMAL(12, 2) NOT NULL,
    status VARCHAR(30) DEFAULT 'created',
    shipping_address TEXT,
    payment_method VARCHAR(50),
    created_at TIMESTAMP DEFAULT CURRENT_TIMESTAMP,
    updated_at TIMESTAMP DEFAULT CURRENT_TIMESTAMP
);

CREATE TABLE order_items (
    id SERIAL PRIMARY KEY,
    order_id INT REFERENCES orders(id),
    product_id INT NOT NULL,
    product_variant_id INT NOT NULL,
    product_name VARCHAR(255),
    variant_name VARCHAR(255),
    unit_price DECIMAL(12, 2) NOT NULL,
    quantity INT NOT NULL,
    subtotal DECIMAL(12, 2) NOT NULL
);

CREATE TABLE order_status_history (
    id SERIAL PRIMARY KEY,
    order_id INT REFERENCES orders(id),
    old_status VARCHAR(30),
    new_status VARCHAR(30),
    note TEXT,
    created_at TIMESTAMP DEFAULT CURRENT_TIMESTAMP
);

CREATE INDEX idx_orders_user_id
ON orders(user_id);

CREATE INDEX idx_orders_status
ON orders(status);
\end{lstlisting}

\subsection{Payment Service Database}

\begin{lstlisting}[style=code,language=SQL]
CREATE TABLE payments (
    id SERIAL PRIMARY KEY,
    order_id INT NOT NULL,
    user_id INT NOT NULL,
    amount DECIMAL(12, 2) NOT NULL,
    payment_method VARCHAR(50),
    status VARCHAR(20) DEFAULT 'pending',
    transaction_code VARCHAR(100),
    payment_url TEXT,
    gateway_response JSONB,
    created_at TIMESTAMP DEFAULT CURRENT_TIMESTAMP,
    updated_at TIMESTAMP DEFAULT CURRENT_TIMESTAMP
);

CREATE INDEX idx_payments_order_id
ON payments(order_id);

CREATE INDEX idx_payments_transaction_code
ON payments(transaction_code);
\end{lstlisting}

\subsection{Shipping Service Database}

\begin{lstlisting}[style=code,language=SQL]
CREATE TABLE shipments (
    id SERIAL PRIMARY KEY,
    order_id INT NOT NULL,
    user_id INT NOT NULL,
    receiver_name VARCHAR(255),
    receiver_phone VARCHAR(20),
    address TEXT NOT NULL,
    carrier VARCHAR(100),
    tracking_code VARCHAR(100),
    status VARCHAR(30) DEFAULT 'processing',
    created_at TIMESTAMP DEFAULT CURRENT_TIMESTAMP,
    updated_at TIMESTAMP DEFAULT CURRENT_TIMESTAMP
);

CREATE INDEX idx_shipments_order_id
ON shipments(order_id);

CREATE INDEX idx_shipments_tracking_code
ON shipments(tracking_code);
\end{lstlisting}

\section{Error Handling}

Để các service trả lỗi thống nhất, hệ thống cần định nghĩa một số mã lỗi chung.

\begin{longtable}{|p{4cm}|p{9cm}|}
\hline
\textbf{Mã lỗi} & \textbf{Ý nghĩa} \\
\hline
UNAUTHORIZED & Người dùng chưa đăng nhập hoặc token không hợp lệ. \\
\hline
FORBIDDEN & Người dùng không có quyền thực hiện thao tác. \\
\hline
PRODUCT\_NOT\_FOUND & Không tìm thấy sản phẩm hoặc biến thể sản phẩm. \\
\hline
OUT\_OF\_STOCK & Sản phẩm không đủ tồn kho. \\
\hline
CART\_EMPTY & Giỏ hàng rỗng, không thể checkout. \\
\hline
ORDER\_NOT\_FOUND & Không tìm thấy đơn hàng. \\
\hline
PAYMENT\_FAILED & Thanh toán thất bại. \\
\hline
INVALID\_ORDER\_STATUS & Trạng thái đơn hàng không hợp lệ. \\
\hline
SERVICE\_UNAVAILABLE & Service phụ thuộc đang lỗi hoặc không phản hồi. \\
\hline
\end{longtable}

Response lỗi nên có cấu trúc thống nhất:

\begin{lstlisting}[style=code]
{
  "error_code": "OUT_OF_STOCK",
  "message": "Product variant does not have enough stock",
  "details": {
    "product_variant_id": 10,
    "requested_quantity": 3,
    "available_stock": 1
  }
}
\end{lstlisting}

\section{Docker và triển khai ban đầu}

Để triển khai hệ thống ở môi trường phát triển, mỗi service có thể được đóng gói bằng Docker. Docker Compose được sử dụng để chạy nhiều service cùng lúc.

Cấu trúc thư mục đề xuất:

\begin{verbatim}
ecommerce-microservices/
  api-gateway/
  user-service/
  product-service/
  cart-service/
  order-service/
  payment-service/
  shipping-service/
  notification-service/
  docker-compose.yml
  README.md
\end{verbatim}

Ví dụ cấu hình môi trường cho từng service:

\begin{lstlisting}[style=code]
USER_SERVICE_PORT=8001
PRODUCT_SERVICE_PORT=8002
CART_SERVICE_PORT=8003
ORDER_SERVICE_PORT=8004
PAYMENT_SERVICE_PORT=8005
SHIPPING_SERVICE_PORT=8006
NOTIFICATION_SERVICE_PORT=8007
GATEWAY_PORT=8080
\end{lstlisting}

Các bước chạy hệ thống:

\begin{enumerate}
  \item Tạo file \texttt{.env} cho từng service.
  \item Chạy database tương ứng cho từng service.
  \item Chạy migration cho từng service.
  \item Seed dữ liệu mẫu cho User Service và Product Service.
  \item Chạy toàn bộ hệ thống bằng Docker Compose.
  \item Kiểm thử API thông qua API Gateway.
\end{enumerate}

\section{Đánh giá thiết kế}

\subsection{Ưu điểm}

\begin{itemize}
  \item Hệ thống được chia theo nghiệp vụ rõ ràng, phù hợp với DDD.
  \item Mỗi service có trách nhiệm riêng, giúp giảm coupling.
  \item Có thể scale từng service độc lập theo tải thực tế.
  \item Thiết kế có Product Variant/SKU nên phù hợp hơn với ecommerce thật.
  \item Sử dụng \texttt{DECIMAL} cho dữ liệu tiền tệ, tránh lỗi sai số.
  \item Checkout flow có kiểm tra giá, tồn kho, giữ hàng và xử lý thanh toán.
  \item Payment flow hỗ trợ webhook, phù hợp với các cổng thanh toán thực tế.
  \item Có state machine cho Order, Payment và Shipping, giúp kiểm soát trạng thái tốt hơn.
\end{itemize}

\subsection{Nhược điểm}

\begin{itemize}
  \item Số lượng service nhiều hơn Monolithic nên vận hành phức tạp hơn.
  \item Cần xử lý lỗi mạng, timeout và retry giữa các service.
  \item Việc đảm bảo nhất quán dữ liệu giữa Order, Payment và Product khó hơn hệ thống dùng chung database.
  \item Cần đầu tư vào logging, monitoring, tracing và CI/CD.
  \item Nếu hệ thống còn nhỏ, triển khai Microservices có thể gây over-engineering.
\end{itemize}

\subsection{Hướng mở rộng}

Trong tương lai, hệ thống có thể mở rộng thêm các service sau:

\begin{itemize}
  \item \textbf{Inventory Service}: Tách riêng quản lý tồn kho, reserve stock và warehouse.
  \item \textbf{Review Service}: Cho phép người dùng đánh giá sản phẩm.
  \item \textbf{Search Service}: Sử dụng Elasticsearch hoặc OpenSearch để tìm kiếm sản phẩm nhanh hơn.
  \item \textbf{Promotion Service}: Quản lý mã giảm giá, voucher và chương trình khuyến mại.
  \item \textbf{AI Service}: Gợi ý sản phẩm hoặc chatbot tư vấn mua hàng.
\end{itemize}

\section{Kết luận}

Chương này đã trình bày thiết kế hệ thống ecommerce theo kiến trúc Microservices ở mức có thể sử dụng làm nền tảng triển khai thực tế. Hệ thống được phân rã theo Domain Driven Design thành các service chính gồm User Service, Product Service, Cart Service, Order Service, Payment Service, Shipping Service và Notification Service.

So với thiết kế đơn giản ban đầu, phiên bản này bổ sung các yếu tố quan trọng của ecommerce thực tế như Product Variant/SKU, xử lý tiền bằng Decimal, cơ chế giữ hàng tạm thời, checkout flow do Order Service điều phối, payment webhook, state machine và error handling thống nhất. Những yếu tố này giúp tài liệu không chỉ dừng lại ở mức mô tả kiến trúc mà còn có thể dùng làm tài liệu định hướng khi xây dựng hệ thống.

Tuy nhiên, Microservices cũng làm tăng độ phức tạp trong triển khai và vận hành. Vì vậy, nếu hệ thống được xây dựng trong phạm vi đồ án hoặc MVP, có thể triển khai trước các service cốt lõi, sau đó mở rộng dần các service nâng cao như Inventory, Search, Promotion và AI Service khi nhu cầu thực tế tăng lên.

\clearpage
```


% =======================
% CHAPTER 3
% =======================